% ========================================================
% BÀI TẬP TOÁN 5 - KẾT NỐI TRI THỨC VỚI CUỘC SỐNG
% Bài 3: Ôn tập phân số
% Định dạng: BTTeX Part III - Trắc nghiệm điền số ngắn
% ========================================================

% Câu 1
\begin{ex}
  Rút gọn phân số $\frac{24}{36}$ về phân số tối giản. Tử số của phân số tối giản là bao nhiêu?
  \shortans{2}
  \loigiai{
    Chia cả tử và mẫu cho $12$: $\frac{24 : 12}{36 : 12} = \frac{2}{3}$. Vậy tử số là 2.
  }
\end{ex}

% Câu 2
\begin{ex}
  Rút gọn phân số $\frac{45}{60}$ về phân số tối giản. Mẫu số của phân số tối giản là bao nhiêu?
  \shortans{4}
  \loigiai{
    Chia cả tử và mẫu cho $15$: $\frac{45 : 15}{60 : 15} = \frac{3}{4}$. Vậy mẫu số là 4.
  }
\end{ex}

% Câu 3
\begin{ex}
  Tìm số tự nhiên $x$ để $\frac{x}{8} = \frac{15}{24}$.
  \shortans{5}
  \loigiai{
    Rút gọn $\frac{15}{24} = \frac{5}{8}$. Do đó $x = 5$.
  }
\end{ex}

% Câu 4
\begin{ex}
  Tìm số tự nhiên $y$ để $\frac{7}{y} = \frac{28}{36}$.
  \shortans{9}
  \loigiai{
    Ta có $\frac{28}{36} = \frac{28 : 4}{36 : 4} = \frac{7}{9}$. Do đó $y = 9$.
  }
\end{ex}

% Câu 5
\begin{ex}
  Quy đồng mẫu số hai phân số $\frac{3}{4}$ và $\frac{5}{6}$ với mẫu số chung nhỏ nhất. Mẫu số chung đó là bao nhiêu?
  \shortans{12}
  \loigiai{
    Mẫu số chung nhỏ nhất của 4 và 6 là $12$ (vì 12 chia hết cho cả 4 và 6).
  }
\end{ex}

% Câu 6
\begin{ex}
  Quy đồng mẫu số hai phân số $\frac{2}{5}$ và $\frac{3}{8}$ với mẫu số chung nhỏ nhất là $40$. Tổng hai tử số sau khi quy đồng là bao nhiêu?
  \shortans{31}
  \loigiai{
    $\frac{2}{5} = \frac{16}{40}$; $\frac{3}{8} = \frac{15}{40}$. Tổng hai tử số là $16 + 15 = 31$.
  }
\end{ex}

% Câu 7
\begin{ex}
  Có bao nhiêu phân số có mẫu số bằng $7$ và nhỏ hơn $1$ (tử số lớn hơn $0$)?
  \shortans{6}
  \loigiai{
    Phân số nhỏ hơn 1 khi tử số nhỏ hơn mẫu số. Các tử số có thể là $1, 2, 3, 4, 5, 6$. Vậy có đúng 6 phân số.
  }
\end{ex}

% Câu 8
\begin{ex}
  Trong các phân số $\frac{5}{6}, \frac{7}{8}, \frac{3}{4}, \frac{11}{12}$, phân số lớn nhất có tử số là bao nhiêu?
  \shortans{11}
  \loigiai{
    Xét phần bù tới 1: $1 - \frac{11}{12} = \frac{1}{12}$ nhỏ nhất nên phân số $\frac{11}{12}$ lớn nhất. Tử số của nó là 11.
  }
\end{ex}

% Câu 9
\begin{ex}
  Trong các phân số $\frac{4}{5}, \frac{2}{3}, \frac{7}{10}, \frac{5}{6}$, phân số bé nhất có mẫu số là bao nhiêu?
  \shortans{3}
  \loigiai{
    Quy đồng mẫu số 30: $\frac{24}{30}, \frac{20}{30}, \frac{21}{30}, \frac{25}{30}$. Phân số bé nhất là $\frac{20}{30} = \frac{2}{3}$. Mẫu số là 3.
  }
\end{ex}

% Câu 10
\begin{ex}
  Cho phân số $\frac{13}{17}$. Cần cộng thêm vào cả tử số và mẫu số của phân số này cùng một số tự nhiên nào để được phân số mới bằng $\frac{4}{5}$?
  \shortans{3}
  \loigiai{
    Hiệu giữa mẫu số và tử số không đổi: $17 - 13 = 4$. Khi phân số mới bằng $\frac{4}{5}$, hiệu giữa mẫu và tử là $5 - 4 = 1$ phần ứng với 4. Tử số mới là $4 \times 4 = 16$. Số cần thêm vào là $16 - 13 = 3$.
  }
\end{ex}

% Câu 11
\begin{ex}
  Tìm số tự nhiên $a$ biết rằng $\frac{2}{3} < \frac{a}{12} < \frac{5}{6}$.
  \shortans{9}
  \loigiai{
    Quy đồng mẫu số 12: $\frac{8}{12} < \frac{a}{12} < \frac{10}{12}$. Suy ra $8 < a < 10 \implies a = 9$.
  }
\end{ex}

% Câu 12
\begin{ex}
  Có bao nhiêu phân số bằng phân số $\frac{3}{5}$ mà tử số là số có hai chữ số?
  \shortans{30}
  \loigiai{
    Tử số có dạng $3 \times k$ với $10 \le 3k \le 99 \implies 3{,}33 \le k \le 33$. Các giá trị nguyên của $k$ từ 4 đến 33. Số lượng phân số là $33 - 4 + 1 = 30$.
  }
\end{ex}

% Câu 13
\begin{ex}
  Tính tổng của tử số và mẫu số của phân số tối giản bằng phân số $\frac{84}{108}$.
  \shortans{16}
  \loigiai{
    Chia cả tử và mẫu cho $12$: $\frac{84 : 12}{108 : 12} = \frac{7}{9}$. Phân số tối giản là $\frac{7}{9}$. Tổng tử số và mẫu số là $7 + 9 = 16$.
  }
\end{ex}

% Câu 14
\begin{ex}
  Một lớp học có $35$ học sinh, trong đó số học sinh nữ chiếm $\frac{3}{5}$ số học sinh cả lớp. Hỏi lớp học đó có bao nhiêu học sinh nam?
  \shortans{14}
  \loigiai{
    Số học sinh nữ là $35 \times \frac{3}{5} = 21$ học sinh. Số học sinh nam là $35 - 21 = 14$ học sinh.
  }
\end{ex}

% Câu 15
\begin{ex}
  Một rổ cam có $48$ quả. Người ta đã lấy đi $\frac{3}{8}$ số cam. Hỏi trong rổ còn lại bao nhiêu quả cam?
  \shortans{30}
  \loigiai{
    Số cam đã lấy đi là $48 \times \frac{3}{8} = 18$ quả. Số cam còn lại là $48 - 18 = 30$ quả.
  }
\end{ex}

% Câu 16
\begin{ex}
  Tìm phân số tối giản $\frac{a}{b}$ biết tổng tử số và mẫu số bằng $25$ và khi rút gọn phân số đó ta được $\frac{2}{3}$. Tử số $a$ là bao nhiêu?
  \shortans{10}
  \loigiai{
    Tổng số phần bằng nhau là $2 + 3 = 5$ phần. Giá trị một phần là $25 : 5 = 5$. Tử số là $5 \times 2 = 10$.
  }
\end{ex}

% Câu 17
\begin{ex}
  Tìm phân số tối giản $\frac{a}{b}$ biết mẫu số hơn tử số $18$ đơn vị và phân số đó bằng $\frac{5}{7}$. Mẫu số $b$ là bao nhiêu?
  \shortans{63}
  \loigiai{
    Hiệu số phần là $7 - 5 = 2$ phần, ứng với 18. Giá trị một phần là $18 : 2 = 9$. Mẫu số $b$ là $9 \times 7 = 63$.
  }
\end{ex}

% Câu 18
\begin{ex}
  Có bao nhiêu phân số lớn hơn $\frac{1}{2}$ nhưng nhỏ hơn $\frac{3}{4}$ có mẫu số bằng $12$?
  \shortans{2}
  \loigiai{
    Quy đồng mẫu 12: $\frac{6}{12} < \frac{x}{12} < \frac{9}{12}$. Tử số $x \in \{7, 8\}$. Vậy có đúng 2 phân số.
  }
\end{ex}

% Câu 19
\begin{ex}
  Tìm số tự nhiên $m$ sao cho $\frac{25 - m}{39} = \frac{6}{13}$.
  \shortans{7}
  \loigiai{
    Ta có $\frac{6}{13} = \frac{6 \times 3}{13 \times 3} = \frac{18}{39}$. Suy ra $25 - m = 18 \implies m = 7$.
  }
\end{ex}

% Câu 20
\begin{ex}
  Phân số $\frac{7}{8}$ gấp phân số $\frac{7}{24}$ bao nhiêu lần?
  \shortans{3}
  \loigiai{
    Thực hiện phép chia: $\frac{7}{8} : \frac{7}{24} = \frac{7}{8} \times \frac{24}{7} = 3$ (lần).
  }
\end{ex}

% Câu 21
\begin{ex}
  Một hình chữ nhật có chiều dài $\frac{4}{5}\text{ m}$, chiều rộng kém chiều dài $\frac{1}{5}\text{ m}$. Chu vi hình chữ nhật đó là bao nhiêu phần mười mét? (Điền tử số khi mẫu số là 10).
  \shortans{28}
  \loigiai{
    Chiều rộng là $\frac{4}{5} - \frac{1}{5} = \frac{3}{5}\text{ m}$. Chu vi là $(\frac{4}{5} + \frac{3}{5}) \times 2 = \frac{14}{5}\text{ m} = \frac{28}{10}\text{ m}$. Tử số là 28.
  }
\end{ex}

% Câu 22
\begin{ex}
  Tìm số tự nhiên $x$ để phân số $\frac{x}{2024} = 0$.
  \shortans{0}
  \loigiai{
    Phân số bằng 0 khi và chỉ khi tử số bằng 0. Do đó $x = 0$.
  }
\end{ex}

% Câu 23
\begin{ex}
  Cho dãy các phân số: $\frac{1}{2}, \frac{1}{4}, \frac{1}{8}, \frac{1}{16}, \dots$. Mẫu số của phân số thứ $6$ trong dãy là bao nhiêu?
  \shortans{64}
  \loigiai{
    Mẫu số của các phân số lần lượt là $2^1, 2^2, 2^3, 2^4, \dots$. Mẫu số thứ 6 là $2^6 = 64$.
  }
\end{ex}

% Câu 24
\begin{ex}
  Tổng của tử số và mẫu số của phân số $\frac{19}{25}$ là bao nhiêu?
  \shortans{44}
  \loigiai{
    Tổng tử số và mẫu số là $19 + 25 = 44$.
  }
\end{ex}

% Câu 25
\begin{ex}
  Một người đi bộ được $\frac{2}{5}$ quãng đường trong giờ thứ nhất, giờ thứ hai đi được $\frac{1}{3}$ quãng đường. Cả hai giờ người đó đi được bao nhiêu phần mười lăm quãng đường? (Điền tử số).
  \shortans{11}
  \loigiai{
    Cả hai giờ đi được: $\frac{2}{5} + \frac{1}{3} = \frac{6}{15} + \frac{5}{15} = \frac{11}{15}$ quãng đường. Tử số là 11.
  }
\end{ex}
